\documentclass[a4paper,12pt]{article}
\usepackage[utf8]{inputenc}
\usepackage[T1]{fontenc}
\usepackage{amsmath,amssymb,amsfonts}
\usepackage{graphicx}
\usepackage{hyperref}
\usepackage{geometry}
\geometry{margin=1in}

\title{Video Processing and Optimisation}
\author{Shashwat Saxena}
\date{\today}

\begin{document}

\maketitle



\section{Introduction}

\subsection{Problem Statement}
Conventional methods for video compression are based on slicing frames and often lead to temporal discontinuty and poor results in video with motion. This research aims to develop a motion-compensated temporal frame fusion algorithem that reduces frames by integrating information from multiple consecutive frames onto one net frame, thereby preserving information and allowing data reduction.

\section{Mathematical Background}
Video could be thought of a 3D signal, with the domain being time (\mathbf{R}) and the co-domain being images (\mathbf{R}). At each instance of time, an image is being displayed, an image itself is a signal that maps a 2D space (\mathbf{R^2}) onto light intensity of 3 colors (RGB), (\mathbf{R x R x R})

\section{Notes}
% Add project notes and equations here.

\end{document}
